% Template for BU ECE RPC written report
\documentclass[12pt,onecolumn,letterpaper]{IEEEtran}  % As per the instructions, use 12pt Times New Roman font
\usepackage[margin=1in]{geometry}  % As per the instructions, use 1-inch margins
\usepackage{titling}  % For customizing the title page
\usepackage{hyperref}  % For fancy hyperlinks in the document
\usepackage{amsmath, amsfonts, amssymb}  % For math stuff
\usepackage{algorithm}  % For algorithms
\usepackage{algpseudocode}  % For algorithmic environment
\usepackage{graphicx}  % For including images
\usepackage[caption=false,font=footnotesize]{subfig}  % For pretty subfigures
\usepackage{mwe}  % Only used for placeholder images, feel free to remove this package
\usepackage{lipsum}  % Only used for placeholder text, feel free to remove this package


% Define theorem-like environments
\newtheorem{definition}{Definition}
\newtheorem{example}{Example}
\newtheorem{proposition}{Proposition}
\newtheorem{remark}{Remark}
\newtheorem{theorem}{Theorem}

% Define customized math operators
\DeclareMathOperator*{\argmax}{arg\,max}
\DeclareMathOperator*{\argmin}{arg\,min}


% Begin
\begin{document}

% Title page
\begin{titlepage}
    \centering
    
    % Title
    {\textbf{TITLE (not identical to paper title)}}
    
    \vspace{5\baselineskip}
    
    % Student name
    {Student name}
    
    \vspace{5\baselineskip}
    
    % Advisor
    {Advisor}
    
    \vspace{5\baselineskip}

    % Abstract
    \begin{abstract}
        \lipsum[1]
    \end{abstract}
    
    \vfill
    
    % Full citation of the main paper, IEEE citation style, 
    \begin{minipage}{\textwidth}
        \footnotesize
        Full citation of the main paper, IEEE citation style. 
    \end{minipage}
    
\end{titlepage}


% Section I
\section{Introduction}
The following is an example of an equation.
\begin{equation}
    E = mc^2
\end{equation}
The following is an example of a multi-line equation.
\begin{align}
    f(x) & = ax^2 + bx + c \\
    & = a(x - r_1)(x - r_2)
\end{align}
You can refer to algorithms (see Algorithm~\ref{alg1}), definitions (see Definition~\ref{def1}), examples (see Example~\ref{eg1}), propositions (see Proposition~\ref{prop1}), remarks (see Remark~\ref{rem1}), theorems (see Theorem~\ref{thm1}), figures (see Fig.~\ref{fig1}), subfigures (see Fig.~\ref{fig1a}, \ref{fig1b}), and tables (see Table~\ref{tab1}). 
You can cite papers from your references \cite{boyd_distributed_2010}.

You can insert links, e.g., \href{https://www.google.com}{https://www.google.com}.
You can also insert footnotes\footnote{This is an example of a footnote.}.


% Section II
\section{Method}
\lipsum[3-5]


% Section III
\section{Critique}
\lipsum[6-8]


% Section IV
\section{Future Work}
\lipsum[9-12]


% Appendix
\newpage
\appendices
\section{Figures and Tables}

% Algorithm
\begin{algorithm}
    \caption{This is an example of an algorithm.} \label{alg1}
    \begin{algorithmic}[1]
        \Procedure{Max}{$a, b$}
            \If{$a > b$}
                \State \Return $a$
            \Else
                \State \Return $b$
            \EndIf
        \EndProcedure
    \end{algorithmic}
\end{algorithm}

% Definition
\begin{definition}[A definition] \label{def1}
    This is an example of a definition.
\end{definition}

% Example
\begin{example}[An example] \label{eg1}
    This is an example of an example.
\end{example}

% Proposition
\begin{proposition}[A proposition] \label{prop1}
    This is an example of a proposition.
\end{proposition}

% Remark
\begin{remark}[A remark] \label{rem1}
    This is an example of a remark.
\end{remark}

% Theorem
\begin{theorem}[A theorem] \label{thm1}
    This is an example of a theorem.
\end{theorem}

% Proof
\begin{IEEEproof}
    This is an example of a proof.
\end{IEEEproof}

% Figure
\begin{figure}[ht!]
    \centering
    \subfloat[]{\includegraphics[width=0.45\linewidth]{example-image}\label{fig1a}}
    \hfil
    \subfloat[]{\includegraphics[width=0.45\linewidth]{example-image}\label{fig1b}}
    \caption{This is an example of a figure. (a) This is the first subfigure. (b) This is the second subfigure.}
    \label{fig1}
\end{figure}

% Table
\begin{table}[ht!]
    \centering
    \caption{This is an example of a table.}
    \begin{tabular}{ccc}
        \hline
        $a$ & $b$ & $\max\{a, b\}$ \\
        \hline
        0 & 0 & 0 \\
        0 & 1 & 1 \\
        1 & 1 & 1 \\
        \hline
    \end{tabular}
    \label{tab1}
\end{table}


% Reference
\newpage
\bibliographystyle{IEEEtran}
\bibliography{refs}  % make sure your references are in `refs.bib`

\end{document}
